\documentclass[12pt]{article}
\usepackage{../style_files/aas_macros}
\usepackage{natbib}
\usepackage[utf8]{inputenc}
\usepackage[left=3cm,right=3cm,top=3cm,bottom=3cm]{geometry}
\usepackage{color}
\usepackage[colorlinks = true,
            linkcolor = blue,
            urlcolor  = blue,
            citecolor = blue]{hyperref}
 \usepackage{graphicx}
 \usepackage{subfigure}
 \usepackage{pythonhighlight}
 \usepackage{comment}
 \usepackage{tablefootnote}
\usepackage{longtable}
\usepackage{tabularx}
\usepackage{booktabs} % toprule, etc.
\usepackage{makecell}
\usepackage{soul}
\usepackage{dirtree}
\sethlcolor{lightgray}
\usepackage{listings}
\lstset{basicstyle=\ttfamily,
  showstringspaces=false,
  commentstyle=\color{red},
  keywordstyle=\color{blue},
  backgroundcolor=\color{lightgray},
  numbers=left
}
\lstdefinestyle{yaml}{
     basicstyle=\color{blue}\footnotesize,
     rulecolor=\color{black},
     string=[s]{'}{'},
     stringstyle=\color{blue},
     comment=[l]{:},
     commentstyle=\color{black},
     morecomment=[l]{-}
}
\usepackage[most]{tcolorbox}

%textmarker style from colorbox doc
\tcbset{textmarker/.style={%
        enhanced,
        parbox=false,boxrule=0mm,boxsep=0mm,arc=0mm,
        outer arc=0mm,left=6mm,right=3mm,top=7pt,bottom=7pt,
        toptitle=1mm,bottomtitle=1mm,oversize}}


% define new colorboxes
\newtcolorbox{hintBox}{textmarker,
    borderline west={6pt}{0pt}{yellow},
    colback=yellow!10!white}
\newtcolorbox{importantBox}{textmarker,
    borderline west={6pt}{0pt}{red},
    colback=red!10!white}
\newtcolorbox{noteBox}{textmarker,
    borderline west={6pt}{0pt}{green},
    colback=green!10!white}

% define commands for easy access
\newcommand{\note}[1]{\begin{noteBox} \textbf{Note:} #1 \end{noteBox}}
\newcommand{\warning}[1]{\begin{hintBox} \textbf{Warning:} #1 \end{hintBox}}
\newcommand{\important}[1]{\begin{importantBox} \textbf{Important:} #1 \end{importantBox}}
\newcommand{\todo}[1]{\begin{hintBox} \textbf{TODO:} #1 \end{hintBox}}

\hypersetup{citecolor=blue}

\definecolor{MSBlue}{rgb}{.204,.353,.541}
\usepackage{titlesec}
\titleformat*{\section}{\color{MSBlue}}
\titleformat*{\subsection}{\color{MSBlue}}
\setlength{\parindent}{0pt}


\title{Residual and ensemble regression for improved reconstruction of high-energy gamma-ray events}
\author{Gernot Maier}
\date{\today}

\begin{document}

\maketitle

\begin{abstract}
We present a machine-learning approach to event reconstruction
for imaging atmospheric Cherenkov telescopes (IACTs) such as VERITAS and
CTAO.
A single XGBoost model acts as a thin layer on top of established methods and
corrects simultaneously energy and direction using multi-target residual regression.
Features are image and stereo parameters based mostly on classical methods like
the Hillas parameters or mean-reduced scaled parameters calculated from MC-trained
lookup tables.
The method is validated on VERITAS data and on simulations of the Cherenkov Telescope
Array Observatory (CTAO) showing a significant improvement in angular and energy
resolution compared to the baseline reconstruction methods.
\end{abstract}

\section{Introduction}

Imaging Atmospheric Cherenkov Telescopes (IACTs) enable Gamma-ray astronomy
to measure the energy and direction of photons from
 tens of GeV to hundreds of TeV (e.g., \citep{2009ARA&A..47..523H}).
The telescopes record Cherenkov light from air showers with
large reflectors, pixelated cameras, fast electronics, and trigger systems.
Modern observatories measure projections of the longitudinal development of air showers
with multiple telescopes simultaneously,
enabling stereoscopic reconstruction of primary direction and energy.
The accuracy of this technique is limited by intrinsic air-shower fluctuations and
telescope photon-detection efficiency to estimated
1 (0.3) arcmin angular resolution at energies of 100 GeV (1 TeV)
\citep{hofmann2006performancelimitscherenkovinstruments}.
Systematic uncertainties from varying observation conditions (e.g., night-sky background light,
atmospheric conditions) further limit the achievable performance and require robust
reconstruction methods.

The different groups operating IACTs have developed a variety of reconstruction methods for direction and energy estimation,
including robust and efficient geometric methods based on image parameters
\citep{1985ICRC....3..445H, 1997APh.....8....1D, 1999APh....12..135H},
template-model based methods \citep{2001A&A...374..895P,2006APh....25..195L, 2009APh....32..231D,2014APh....56...26P},
and machine-learning-based methods using
Hillas parameters \citep{2024A&A...691A.328A} or full camera images.
Performance varies across methods, with template-model based algorithms achieving
the highest precision at the cost of computational complexity and robustness.
\todo{Missing citations for FreePact / full image methods.}

This paper shows that a lightweight machine-learning regression layer
using the XGBoost library \citep{Chen:2016:XST:2939672.2939785}
can improve the performance of existing image-parameter-based
reconstruction methods.
The method implements \textit{residual regression}: rather than predicting direction
and energy directly, a single multi-target XGBoost model learns the corrections
(residuals) to an existing baseline reconstruction, effectively combining the robustness
of the baseline method with the flexibility of gradient-boosted trees.
The explicit goal is to apply these models to archival data spanning almost two
decades of VERITAS \citep{2006APh....25..391H} observations
without requiring a complete re-write of the analysis code or low-level reprocessing
of raw data.
The method is sufficiently general to be easily adapted to different IACT arrays,
as demonstrated for both VERITAS and simulations of the Northern array of the Cherenkov Telescope Array Observatory (CTAO)
\citep{2019scta.book.....C}.

The code developed for this paper is available from \citep{maier_2026_18412988}.
All methods and results presented here are based on the open-source Eventdisplay reconstruction
code \citep{maier_2025_17447664,maier_2026_18412887}.

The paper is structured as follows:
Section~\ref{sec:features_and_targets} describes the data pipeline, feature engineering,
and regression targets.
Section~\ref{sec:regression} presents the regression model architecture and training strategy.
Section~\ref{sec:performance} evaluates performance on simulated events,
and Section~\ref{sec:results} presents results on Crab Nebula data.
Section~\ref{sec:systematics} discusses systematic effects and robustness.
Section~\ref{sec:conclusions} concludes with a summary and outlook.


\section{Data Pipeline, Features, and Regression Targets}
\label{sec:features_and_targets}

\subsection{Input data and pre-processing}

The analysis uses the output of the Eventdisplay stereo reconstruction,
stored in ROOT-format mscw files.
These files contain, for each triggered event, the
Hillas image parameters for each telescope, the results of two
existing reconstruction methods (see below), and the Monte Carlo (MC) truth values
used for training (primary direction offsets and energy).
Pre-selection requires a minimum of two images passing image cleaning
(\texttt{DispNImages} $\geq 2$).
No additional quality cuts are applied at the training stage to keep the
method applicable to the full range of observation conditions.

Variables with large dynamic range are log$_{10}$-transformed before use
(image size, energy estimates, pixel counts, core distance).
Physical limits and robustness against badly reconstructed events are
enforced through clipping; for example, energy estimates are clipped to
$E > 10^{-3}$~TeV, image sizes to $> 1$~p.e., and field-of-view offsets
to $|\Delta\theta| < 7$~deg.
Missing telescope data (telescopes not participating in a given event)
are filled with \texttt{NaN} values; the XGBoost
histogram-based tree algorithm (\texttt{tree\_method=hist}) handles
these natively without explicit imputation.

\subsection{Feature engineering and sorting}
\label{sec:features}

Features are organised in three groups: telescope-level image features,
stereo-reconstruction features, and event-level array and observation features.
A detailed listing is provided in Appendix~\ref{sec:appendix_features}.

\textbf{Telescope-level features} are derived from image (Hillas) parameterisation
of each camera image: major and minor axis lengths (\texttt{length}, \texttt{width}),
image orientation (\texttt{cosphi}, \texttt{sinphi}), image centroid
(\texttt{cen\_x}, \texttt{cen\_y}), image size (log$_{10}$, \texttt{size}),
image asymmetry (\texttt{asym}), time gradient along the image axis (\texttt{tgrad\_x}),
and fractional image loss at the camera edge (\texttt{loss}).
Additional telescope-level features are
the per-telescope energy estimates from the DispBDT reconstruction
(\texttt{E}, \texttt{ES}), the number of image pixels after cleaning (\texttt{ntubes}),
the number of high-charge pixels recorded in the low-gain channel (\texttt{nlowgain}),
the impact distance from the reconstructed shower core (\texttt{R\_core}),
and the telescope position relative to the shower core projected onto the shower plane
(\texttt{tel\_rel\_x}, \texttt{tel\_rel\_y}).

To apply a single model to events with different telescope multiplicities and
array configurations, telescope-level features are sorted by image size
(largest image first); telescopes without a valid image are placed last and
their features are set to \texttt{NaN}.
For heterogeneous arrays (e.g., CTAO), the sorting is hierarchical: first by
mirror area (largest telescope type first), then by image size within each type.
This scheme preserves the physics information of the brightest, most constraining
images while allowing the same model to process all multiplicities.

\textbf{Stereo-reconstruction features} include the results of two independent
direction and energy estimators available in Eventdisplay:
(i) the DispBDT reconstruction, which regresses the displacement of the shower
direction from each image centroid; it provides weighted averages of the
predicted pointing offset (\texttt{Xoff}, \texttt{Yoff}) and energy
(\texttt{ErecS}), as well as per-telescope estimates and uncertainties;
and (ii) the intersection-of-lines method, which intersects the image major axes
from all telescope pairs to estimate the shower direction (\texttt{Xoff\_intersect},
\texttt{Yoff\_intersect}).
The DispBDT per-telescope quantities (\texttt{DispXoff\_T}, \texttt{DispYoff\_T},
\texttt{Disp\_T}, \texttt{DispWoff\_T}) carry information on the consistency
and precision of the individual-image reconstruction.
Further stereo features are the Cherenkov emission height (\texttt{EmissionHeight}),
the reconstructed shower core position (\texttt{Xcore}, \texttt{Ycore}),
and the telescope pointing direction (elevation, azimuth).

\textbf{Event-level features} capture the observational geometry and array
configuration: the number of images used by the DispBDT reconstruction
(\texttt{DispNImages}), the 2D convex-hull footprint area of the active telescopes,
the angle between the shower direction and the geomagnetic field
(\texttt{Geomagnetic\_Angle}), and per-telescope mirror areas and activity flags.
The geomagnetic angle is computed from the pointing direction and local field
components specific to each observatory (VERITAS, CTAO-North, CTAO-South),
and accounts for asymmetric deflection of shower secondaries that affects
image shape and orientation.

In total, the feature vector for each event contains more than 100 features
for a four-telescope array.

\subsection{Regression targets}
\label{sec:targets}

The regression model simultaneously predicts three quantities:
the residuals of the DispBDT direction estimate in camera coordinates,
$\Delta x = x_\mathrm{MC} - x_\mathrm{DispBDT}$ (\texttt{Xoff\_residual}) and
$\Delta y = y_\mathrm{MC} - y_\mathrm{DispBDT}$ (\texttt{Yoff\_residual}),
and the residual of the DispBDT energy estimate in log$_{10}$ space,
$\Delta\log E = \log_{10}(E_\mathrm{MC}) - \log_{10}(E_\mathrm{DispBDT})$ (\texttt{E\_residual}).
During inference, the final predictions are obtained by adding the predicted
residuals back to the DispBDT baseline:
\begin{align}
  x_\mathrm{pred} &= x_\mathrm{DispBDT} + \Delta\hat{x}, \\
  y_\mathrm{pred} &= y_\mathrm{DispBDT} + \Delta\hat{y}, \\
  \log_{10}(\hat{E}) &= \log_{10}(E_\mathrm{DispBDT}) + \Delta\hat{\log E}.
\end{align}
Predicting residuals rather than absolute values reduces the learning task
to a correction problem: the model captures systematic biases and conditions
under which the baseline method performs poorly
(e.g., edge-of-camera images, high night-sky background, specific geomagnetic orientations),
while inheriting the robustness of the DispBDT reconstruction for typical events.

\subsection{Performance metrics}
\label{sec:metrics}

Angular performance is quantified by the angular resolution $R_{68}$,
defined as the 68\% containment radius of the distribution of
$\Delta\theta = \sqrt{(\hat{x} - x_\mathrm{MC})^2 + (\hat{y} - y_\mathrm{MC})^2}$
in bins of true energy.
Energy performance is quantified by the 68\% containment of the relative
energy error $\Delta E / E = |10^{\hat{\log E}} - E_\mathrm{MC}| / E_\mathrm{MC}$
(energy resolution) and the mean relative bias
$(10^{\hat{\log E}} - E_\mathrm{MC}) / E_\mathrm{MC}$ (energy bias).
Both metrics are computed in nine bins of $\log_{10}(E/\mathrm{TeV})$
spanning $[-2.0, 2.5]$.



\section{Regression Model}
\label{sec:regression}

\subsection{Model architecture}

A single \texttt{XGBRegressor} from the XGBoost library \citep{Chen:2016:XST:2939672.2939785}
is trained on the three targets simultaneously using the
native multi-output regression mode.
The squared-error objective (\texttt{reg:squarederror}) minimises
$\sum_i w_i (\hat{y}_i - y_i)^2$ jointly across all targets.
Before training, each target is standardised to zero mean and unit variance
(computed from the training set only) to prevent the larger numerical scale of
the energy residual from dominating the direction residuals.
Predictions are inverse-transformed to the original scale during inference.

The histogram-based tree algorithm (\texttt{tree\_method=hist}) is used throughout,
which natively handles \texttt{NaN} feature values (missing telescopes)
by assigning them to the learned default direction at each tree split.
The model therefore processes events with any telescope multiplicity
$\geq 2$ without pre-grouping by multiplicity class.

\subsection{Hyperparameters}

The hyperparameters are summarised in Table~\ref{tab:hyperparameters}.
The combination of a small learning rate (0.02) with a large maximum number of
estimators (10\,000) and early stopping (patience 50) allows fine-grained
convergence while preventing overfitting.
Tree depth (7) and minimum child weight (10) control model complexity.
Column and row subsampling (0.7 each) provide implicit regularisation.

\begin{table}[h]
\centering
\caption{XGBoost regression hyperparameters.}
\label{tab:hyperparameters}
\begin{tabular}{ll}
\toprule
\textbf{Parameter} & \textbf{Value} \\
\midrule
\texttt{objective}           & \texttt{reg:squarederror} \\
\texttt{tree\_method}        & \texttt{hist} \\
\texttt{n\_estimators}       & 10\,000 (with early stopping) \\
\texttt{early\_stopping\_rounds} & 50 \\
\texttt{eval\_metric}        & RMSE \\
\texttt{learning\_rate}      & 0.02 \\
\texttt{max\_depth}          & 7 \\
\texttt{min\_child\_weight}  & 10 \\
\texttt{subsample}           & 0.7 \\
\texttt{colsample\_bytree}   & 0.7 \\
\bottomrule
\end{tabular}
\end{table}

\subsection{Training strategy and sample weighting}
\label{sec:training}

Training data consist of Monte Carlo simulated gamma-ray events after the standard
Eventdisplay stereo reconstruction.
Events are split 50/50 into training and test sets using a fixed random state
to ensure reproducibility and avoid data leakage.
The evaluation set used for early stopping consists of both splits;
the final model is selected at the iteration that minimises test RMSE.

The training set is reweighted by a combined event weight
\begin{equation}
  w_i = w_i^{(E)} \cdot w_i^{(N)},
\end{equation}
where $w^{(E)}$ is an inverse-frequency weight in bins of true log$_{10}$ energy
(events in under-populated energy bins receive higher weight)
and $w^{(N)} \propto N_\mathrm{images}^2$ favours high-multiplicity events.
Both weight components are normalised so that their mean equals unity,
preserving the effective learning rate.
Events in energy bins with fewer than 10 entries receive zero weight and are excluded.
This combined weighting scheme improves reconstruction accuracy at high energies
(where simulated event counts are low) and for the highest-multiplicity events
(where reconstruction constraints are strongest).

\subsection{Feature importance and model interpretability}

Feature importance is assessed using SHAP (SHapley Additive exPlanations) values
computed on a random subsample of 1\,000 test events.
For the native multi-output XGBoost model, SHAP values are extracted per target
from the tree contributions (\texttt{pred\_contribs=True}),
giving separate importance rankings for the direction
($\Delta x$, $\Delta y$) and energy ($\Delta\log E$) corrections.
SHAP values are also evaluated in bins of true energy to identify
which features drive the correction at different energy scales.



\section{Performance Evaluation on Simulated Events}
\label{sec:performance}

\subsection{Angular resolution}

\todo{Figure: Angular resolution $R_{68}$ (68\% containment of $\Delta\theta$) vs.\ $\log_{10}(E/\mathrm{TeV})$ for the XGBoost model and the DispBDT baseline. Show also the 90\% containment.  Generated by \texttt{calculate\_resolution()} in \texttt{evaluate.py}.}

\subsection{Energy resolution and bias}

\todo{Figures: (1) Energy resolution (68\% containment of $|\Delta E/E|$) vs.\ $\log_{10}(E/\mathrm{TeV})$, XGBoost vs.\ DispBDT/ErecS baseline. (2) Energy bias (mean of $\Delta E/E$) vs.\ $\log_{10}(E/\mathrm{TeV})$, same comparison. Generated by \texttt{calculate\_resolution()} in \texttt{evaluate.py}.}

\subsection{Diagnostic checks}

Model diagnostics are evaluated with the following set of plots:

\paragraph{Training convergence.}
The RMSE of the training and test sets is recorded at each boosting round.
Convergence and the absence of significant overfitting is verified by
confirming that the test RMSE plateaus before early stopping triggers,
and that the train--test gap remains below 10\%.
\todo{Figure: Training and test RMSE vs.\ boosting round (all three targets on one plot or one panel each). Mark the early-stopping iteration. Generated by \texttt{plot\_training\_evaluation.py}.}

\paragraph{Generalisation gap.}
The per-target RMSE ratio $R_\mathrm{gen} = \mathrm{RMSE}_\mathrm{test} / \mathrm{RMSE}_\mathrm{train}$
is computed for each of the three targets.
Values close to 1 confirm that the model generalises well to unseen events.
\todo{Figure: Bar chart of train and test RMSE per target ($\Delta x$, $\Delta y$, $\Delta\log E$), and the corresponding gap percentage. Generated by \texttt{diagnostic\_generalization\_gap.py}.}

\paragraph{Residual normality.}
For each target the distribution of prediction errors $(y - \hat{y})$ is tested
for Gaussianity using the Kolmogorov--Smirnov and Anderson--Darling tests.
Heavy tails ($|$kurtosis$| > 1$) or significant skewness indicate systematic failures
for specific event classes.
\todo{Figure: Residual histograms with Gaussian fits and Q-Q plots for each of the three targets. Report KS p-value, skewness, excess kurtosis, and 3$\sigma$ outlier fraction. Generated by \texttt{diagnostic\_residual\_normality.py}.}

\paragraph{SHAP feature importance.}
The top-20 features by mean absolute SHAP value are shown for each target
(direction $x$, direction $y$, energy) to identify which input quantities
drive the corrections.
\todo{Figure: Horizontal bar chart of the top-20 SHAP importances for each target (three panels). Generated by \texttt{diagnostic\_shap\_summary.py}.}

\paragraph{SHAP by energy bin.}
SHAP values are computed separately in nine bins of true energy to
identify energy-dependent feature importance.
\todo{Figure/Table: Heat map or table of mean $|$SHAP$|$ for the top features vs.\ energy bin, one panel per target. Generated by \texttt{shap\_feature\_importance\_by\_energy()} in \texttt{evaluate.py}.}



\section{Results on Crab Nebula Data}
\label{sec:results}

The Crab Nebula, the standard candle of TeV gamma-ray astronomy, is used to
validate the reconstruction performance on observational data.
VERITAS observations of the Crab Nebula provide a large sample of high-significance
gamma-ray events at known sky position, allowing direct comparison of reconstructed
directions and the detection significance between the XGBoost and baseline methods.

% TODO: Include the following result plots.
\todo{The following plots compare XGBoost vs.\ DispBDT baseline on actual VERITAS Crab Nebula observations.}

\paragraph{Theta-squared distribution.}
Distribution of $\theta^2$, the squared angular distance between the reconstructed
event direction and the Crab Nebula position (ON region) and equivalent background
control regions (OFF regions), for the XGBoost and DispBDT methods.
An improvement in angular resolution is expected to manifest as a narrower
ON-region excess and improved signal-to-noise ratio.
\todo{Figure: $\theta^2$ distributions (ON and OFF regions) for XGBoost and DispBDT. Quote Li\&Ma significance, gamma-ray excess counts, and observing time.}

\paragraph{Cumulative angular distributions.}
Cumulative distribution of $\theta$ for ON and OFF events, used to directly compare
the 68\% and 80\% containment radii between methods.
\todo{Figure: Cumulative $\theta$ distribution for ON and OFF events, XGBoost vs.\ DispBDT; mark 68\% and 80\% containment radii.}

\paragraph{Significance map.}
Sky map of detection significance in the Crab Nebula field,
comparing the XGBoost-reconstructed events with the DispBDT baseline.
\todo{Figure: Significance sky map (excess map) in the Crab Nebula field, XGBoost vs.\ DispBDT. Show significance contours.}

\paragraph{Energy spectrum.}
Differential energy spectrum of the Crab Nebula estimated from the
XGBoost-reconstructed events, compared to published measurements
\citep{2015ICRC...34..792A} as a cross-check of the energy reconstruction.

\section{Systematic Effects and Robustness}
\label{sec:systematics}

\subsection{Zenith-angle dependence}

Shower morphology and Cherenkov light yield depend strongly on zenith angle,
affecting both image shapes and the accuracy of reconstruction.
The reconstruction performance (angular and energy resolution) is evaluated
in bins of zenith angle to verify that the model generalises across the
range of zenith angles used in VERITAS observations ($\sim 0$--$65^\circ$).
\todo{Figure: Angular resolution $R_{68}$ and energy resolution vs.\ zenith angle (in bins), XGBoost and DispBDT baseline.}

\subsection{Pointing-offset dependence}

Sources are observed at a range of angular distances from the camera centre
(typically $0.5$--$1.5$~deg for VERITAS).
Image truncation and reduced image quality at large offsets degrade reconstruction
accuracy. The \texttt{loss} feature (fractional image truncation) is explicitly
included in the feature set to allow the model to correct for this effect.
Performance is evaluated as a function of pointing offset.
\todo{Figure: Angular and energy resolution vs.\ source offset from camera centre (in bins), XGBoost vs.\ baseline.}

\subsection{Telescope multiplicity}

A single model trained on all multiplicities $\geq 2$ is evaluated separately
for two-, three-, and four-telescope events.
The feature sorting by image size ensures that the model can handle reduced
multiplicities without retraining, but performance is expected to degrade
at lower multiplicities due to the reduced geometric constraints.
\todo{Figure: Angular and energy resolution separately for 2-, 3-, and 4-telescope events; compare XGBoost and DispBDT.}

\subsection{Domain shift}

The model is trained on simulated events and applied to observational data.
Potential domain shifts arise from imperfect simulation of telescope response,
night-sky background, or atmospheric conditions.
The agreement between MC and data distributions for key features is checked,
and the consistency of the reconstructed energy spectrum with published results
provides an overall validation of the energy calibration.
\todo{Figure: Crab Nebula differential energy spectrum (XGBoost reconstruction) overlaid with published VERITAS/MAGIC/HESS measurements \citep{2015ICRC...34..792A}. Quote flux normalisation and spectral index.}



\section{Conclusions}
\label{sec:conclusions}

This paper presented a machine-learning approach to event reconstruction for ground-based IACTs, demonstrated for VERITAS and
CTAO-North simulations.

For direction and energy reconstruction, a single multi-target XGBoost regression model
simultaneously predicts corrections to the existing DispBDT reconstruction
in all three coordinates (two angular and one energy).
The model is trained on more than 100 image, stereo, and array-geometry features,
using sample weights that balance sparse high-energy bins and favour high-multiplicity events.
Predicting residuals rather than absolute values retains the robustness of the
existing baseline for typical events while learning systematic corrections where the
baseline fails.

The code is designed as a lightweight wrapper around the existing Eventdisplay
pipeline, operating on the mscw output files without requiring reprocessing
of raw data.
This makes it directly applicable to the full VERITAS archival data set
spanning almost two decades of observations.
The same interface supports CTAO simulations with minimal reconfiguration.

% TODO: Summarise quantitative improvements in angular resolution, energy resolution,
% and energy bias once result plots are available.
\todo{Fill in quantitative summary once result plots are generated: angular resolution improvement (\%), energy resolution improvement (\%), energy bias reduction, and Li\&Ma significance gain on Crab data.}

Limitations of the current approach include the use of a single global model
(no zenith-angle binning for regression), reliance on DispBDT as the baseline
(systematic errors in DispBDT propagate to the XGBoost predictions),
and training on Monte Carlo only (domain-shift corrections not applied).
Future work will investigate training on real data using self-supervised methods,
zenith-binned models, and the integration of additional observational features
such as night-sky background level and atmospheric monitoring data.



\section{Detailed description of features}
\label{sec:appendix_features}

\subsection{Telescope-level features}

Telescope-level features are indexed 0 to $N_\mathrm{tel}-1$, sorted by descending
image size (largest first). Features for telescopes without a valid image are set to \texttt{NaN}.
For CTAO, sorting is first by mirror area (largest telescope type first), then by image size.

\begin{longtable}{lll p{5.5cm}}
\caption{Telescope-level features used for stereo reconstruction.
Subscripts 0\ldots$N-1$ denote the sorted telescope index.} \label{tab:features-image} \\
\hline
\textbf{Feature} & \textbf{Unit} & \textbf{Range (clipped)} & \textbf{Description} \\
\hline
\endfirsthead
\hline
\textbf{Feature} & \textbf{Unit} & \textbf{Range (clipped)} & \textbf{Description} \\
\hline
\endhead
\hline
\endfoot
\texttt{size} & log$_{10}$ p.e. & $[0,\infty)$ & Image charge (after $\log_{10}$ transform); minimum 1~p.e.\ before transform \\
\texttt{width} & deg & unconstrained & Hillas minor-axis RMS width \\
\texttt{length} & deg & unconstrained & Hillas major-axis RMS length \\
\texttt{asym} & deg & unconstrained & Hillas image asymmetry \\
\texttt{cosphi} & -- & $[-1,1]$ & Cosine of image major-axis orientation angle \\
\texttt{sinphi} & -- & $[-1,1]$ & Sine of image major-axis orientation angle \\
\texttt{cen\_x} & deg & unconstrained & Hillas centroid $x$ in camera coordinates \\
\texttt{cen\_y} & deg & unconstrained & Hillas centroid $y$ in camera coordinates \\
\texttt{dist} & deg & unconstrained & Distance of image centroid from camera centre \\
\texttt{loss} & -- & $[0,1]$ & Fraction of image charge outside camera boundary \\
\texttt{tgrad\_x} & p.e./ns & $[-50,50]$ & Time gradient along image major axis \\
\texttt{ntubes} & log$_{10}$ & $[0,\infty)$ & Number of pixels surviving image cleaning \\
\texttt{nlowgain} & log$_{10}$ & $[0,\infty)$ & Number of pixels in low-gain readout channel \\
\texttt{E} & log$_{10}$ TeV & $[-3,\infty)$ & Per-telescope energy estimate from DispBDT \\
\texttt{ES} & log$_{10}$ TeV & $[-3,\infty)$ & Per-telescope energy from look-up table \\
\texttt{Disp\_T} & deg & unconstrained & DispBDT displacement magnitude per telescope \\
\texttt{DispXoff\_T} & deg & $[-7,7]$ & DispBDT $x$-offset per telescope \\
\texttt{DispYoff\_T} & deg & $[-7,7]$ & DispBDT $y$-offset per telescope \\
\texttt{DispWoff\_T} & deg & unconstrained & DispBDT weighted offset per telescope \\
\texttt{R\_core} & m & $[-10,\infty)$ & Impact distance from shower core to telescope \\
\texttt{mirror\_area} & m$^2$ & $>0$ & Mirror collection area (constant per telescope type) \\
\texttt{tel\_rel\_x} & m & unconstrained & Telescope $x$ position relative to shower core, shower-plane coords \\
\texttt{tel\_rel\_y} & m & unconstrained & Telescope $y$ position relative to shower core, shower-plane coords \\
\texttt{tel\_active} & -- & $\{0,1\}$ & Flag: 1 if telescope participated in this event \\
\end{longtable}

\subsection{Stereo-reconstruction and event-level features}

\begin{longtable}{lll p{5.5cm}}
\caption{Stereo-reconstruction and event-level features.} \label{tab:features-stereo} \\
\hline
\textbf{Feature} & \textbf{Unit} & \textbf{Range (clipped)} & \textbf{Description} \\
\hline
\endfirsthead
\hline
\textbf{Feature} & \textbf{Unit} & \textbf{Range (clipped)} & \textbf{Description} \\
\hline
\endhead
\hline
\endfoot
\texttt{Xoff} & deg & $[-7,7]$ & DispBDT stereo $x$-offset (weighted average) \\
\texttt{Yoff} & deg & $[-7,7]$ & DispBDT stereo $y$-offset (weighted average) \\
\texttt{Xoff\_intersect} & deg & $[-7,7]$ & Intersection-of-lines $x$-offset \\
\texttt{Yoff\_intersect} & deg & $[-7,7]$ & Intersection-of-lines $y$-offset \\
\texttt{ErecS} & log$_{10}$ TeV & $[-3,\infty)$ & DispBDT stereo energy estimate \\
\texttt{Erec} & log$_{10}$ TeV & $[-3,\infty)$ & Look-up table energy estimate \\
\texttt{EmissionHeight} & km & $[0,100]$ & Reconstructed Cherenkov emission height \\
\texttt{Xcore} & m & unconstrained & Shower core $x$ position (ground) \\
\texttt{Ycore} & m & unconstrained & Shower core $y$ position (ground) \\
\texttt{img2\_ang} & deg & $[0,360]$ & Angle between image pairs (stereoscopic angle) \\
\texttt{DispNImages} & -- & $[2, N_\mathrm{tel}]$ & Number of telescopes used by DispBDT \\
\texttt{ArrayPointing\_Elevation} & deg & $[0,90]$ & Array pointing elevation \\
\texttt{ArrayPointing\_Azimuth} & deg & $[0,360]$ & Array pointing azimuth \\
\texttt{array\_foot\_print} & m$^2$ & $[0,\infty)$ & Convex-hull area of active telescopes \\
\texttt{Geomagnetic\_Angle} & deg & $[0,180]$ & Angle between shower direction and geomagnetic field \\
\end{longtable}

\section*{AI usage disclosure}

Generative AI tools (including Claude, ChatGPT, and Gemini) were used to assist
with code development and debugging of the code associated to this paper.
AI tools are also used for language improvement of this text.
All AI-assisted outputs were reviewed, validated, and, where necessary,
modified by the authors to ensure accuracy and reliability.

\section*{Acknowledgments}

Add acknowledgments here.

%%%%%%%%%%%%%%%%%%%%%%%%%%%%%%%%%%%%%%%%%%%
\bibliographystyle{aasjournal}
\bibliography{../bibliography/bibliography.bib}

\end{document}
